\section{Abstract}
{\justify This project aims to develop a precise location system for an augmented reality application that guides users around the campus of the Universidad del Valle de Guatemala. The system is based on Estimote Beacons with Ultra-Wideband (UWB) technology and is designed to integrate with Unity through a native plugin that enables real-time communication between mobile devices and location sensors.

Unlike previous solutions that relied on QR codes or sensor fusion techniques subject to cumulative errors, this proposal employs Kalman filter-assisted trilateration to more reliably estimate the user's position in 2D coordinates. The system will be able to transmit this data to Unity, where it will be used to render augmented reality elements that guide the user around the campus.

The project is part of a larger ongoing initiative and focuses on the technical implementation of the positioning module, laying the groundwork for future cross-platform extensions. This technological solution significantly contributes to improving the on-campus wayfinding experience and represents an example of the practical use of emerging technologies in localization, signal filtering, and interactive mobile application development.}
\newpage

\section{Resumen}
{\justify
Este proyecto tiene como objetivo desarrollar un sistema de localización precisa para una aplicación de realidad aumentada que guía a los usuarios dentro del campus de la Universidad del Valle de Guatemala. El sistema se basa en sensores Estimote Beacons con tecnología Ultra-Wideband (UWB) y está diseñado para integrarse con Unity mediante un plugin nativo que permite la comunicación en tiempo real entre los dispositivos móviles y los sensores de localización.

A diferencia de soluciones anteriores que dependían de códigos QR o técnicas de sensor fusion sujetas a errores acumulativos, esta propuesta emplea trilateración asistida por filtros de Kalman para estimar con mayor estabilidad la posición del usuario en coordenadas 2D. El sistema será capaz de transmitir estos datos a Unity, donde se utilizan para renderizar elementos de realidad aumentada que guían al usuario por el campus.

El proyecto forma parte de una iniciativa mayor en curso y se enfoca en la implementación técnica del módulo de posicionamiento, sentando las bases para futuras extensiones multiplataforma. Esta solución tecnológica contribuye significativamente a mejorar la experiencia de orientación dentro del campus y representa un ejemplo del uso práctico de tecnologías emergentes en localización, filtrado de señales y desarrollo de aplicaciones móviles interactivas.
}
\newpage